\scanpage{11}

\mspage{5}

On aura une réciproque partielle dans les catégories accessibles. C'est pour l'énoncé de cette réciproque que j'ai mis à part la condition b) de stabilité par facteurs directs, au lieu de la bloquer avec a).

\uline{Définition 1} \quad a) (Pour mémoire) RLP, LLP, la relation $\Phi\longleftrightarrow\Psi$ entre parties de $\mathrm{Fl}\,M$, (relation non symétrique !).

b) (Pour mémoire) Parties de $\mathrm{Fl}\,M$ closes à gauche, \unsure{close|closes} \unsure{à} droite.

c) Soit $\Phi\longleftrightarrow\Psi$. On dit que le couple $(\Phi,\Psi)$ a la propriété de factorisation de Quillen, si \uline{\unsure{toute|tte}} flèche $f$ de $M$ se factorise en

\[ f=pi,\qquad i\in\Phi,\ p\deleted{\unsure{\Psi}}\in\Psi. \]

On dit que $(\Phi,\Psi)$ est un \uline{couple de Quillen}, si \deleted{\unsure{$\Phi$}} c'est un couple associé par propriétés LLP et RLP $(\Phi\unsure{\leftarrow|\leftarrow\!\shortmid}\Psi)$, \deleted{si elles} et elles sont closes (i.e. $\Phi=\Psi^*,\ \Psi=\Phi_*$), et si elles ont la propriété de factorisation de Quillen\unsure{!|.}

J'énonce dès maintenant le résultat technique essentiel, assurant l'existence de beaucoup de couples de Quillen :

% archive sheet number: 10

